\section*{Objetivo}
\justifying

% Verbo infinitivo + ¿Qué? + ¿Para qué? + ¿mediante qué?
-- No se si hay que hacer más de un objetivo, o clasificarlos en general y específicos. --

Diseñar un sistema didáctico modular para la experimentación en electrónica de potencia,
mediante el uso de módulos especializados y material didáctico teórico.

El objetivo general surge directamente del proyecto a estudiar. Es precisamente el “qué”
se va a ofrecer al término del mismo, de aquí que define también sus alcances. En 
el proceso del proyecto, es tan importante la función del objetivo, que si se carece
de él o su redacción no es clara, no existirá una referencia que indique al alumno
si logró lo deseado.

Para plantearlo, ayudaría responder reflexivamente a la pregunta: ¿cuál es la finalidad
del proyecto? El enunciado debe ser claro y preciso; será mejor en cuanto no de a lugar
a múltiples interpretaciones. Debe evitarse englobar todos los objetivos de la
investigación en un solo enunciado.

En resumen, en cuanto a tipo de objetivos podemos señalar lo siguiente: el objetivo
general es la meta que se pretende alcanzar. Define los alcances del estudio y del proyecto.
